\documentclass[10pt]{beamer}
\usetheme{Madrid}
\usepackage[T1]{fontenc}
\usepackage[utf8]{inputenc}
\usepackage[brazil]{babel}
\usepackage{lmodern}
\usepackage{hyperref}
\usepackage{listings}
\usepackage{xcolor}
\setbeamertemplate{navigation symbols}{}

\lstset{
  language=C++,
  basicstyle=\ttfamily\small,
  keywordstyle=\color{blue},
  commentstyle=\color{gray},
  stringstyle=\color{red},
  breaklines=true,
  showstringspaces=false,
}

\title{Arrays e Laços}
\subtitle{Curso de Microcontroladores}
\author{}
\date{}

\begin{document}

\begin{frame}
  \titlepage
\end{frame}

\begin{frame}[fragile]{O que é um array?}
\begin{lstlisting}
int sensores[5];  // Array de 5 inteiros
sensores[0] = 100;
sensores[1] = 200;
sensores[2] = 150;
// ... e assim por diante

for (int i = 0; i < 5; i++) {
  Serial.println(sensores[i]);
}
\end{lstlisting}

  \begin{itemize}
    \item Array é uma lista de valores do mesmo tipo.
    \item Índice começa em 0.
    \item Declarar: `tipo nome[tamanho];`
  \end{itemize}
\end{frame}

\begin{frame}[fragile]{Arrays de sensores do seguidor}
\begin{lstlisting}
const int NUM_SENSORES = 5;
int pinsSensores[NUM_SENSORES] = {
  A0, A1, A2, A3, A4
};
int leituras[NUM_SENSORES];

void lerSensores() {
  for (int i = 0; i < NUM_SENSORES; i++) {
    leituras[i] = analogRead(pinsSensores[i]);
    Serial.print(leituras[i]);
    Serial.print(" ");
  }
  Serial.println();
}
\end{lstlisting}
\end{frame}

\begin{frame}{Tipos de laço}
  \begin{itemize}
    \item `for`: quando sabe quantas vezes repetir.
    \item `while`: enquanto uma condição for verdadeira.
    \item `do...while`: executa pelo menos uma vez.
  \end{itemize}
\end{frame}

\begin{frame}[fragile]{Laço for}
\begin{lstlisting}
// Forma clássica
for (int i = 0; i < 10; i++) {
  Serial.println(i);
}

// Forma com array
int valores[] = {10, 20, 30, 40, 50};
for (int i = 0; i < 5; i++) {
  Serial.println(valores[i]);
}
\end{lstlisting}

  \begin{itemize}
    \item Inicialização, condição, incremento.
    \item Ótimo para percorrer arrays.
  \end{itemize}
\end{frame}

\begin{frame}[fragile]{Laço while}
\begin{lstlisting}
int contador = 0;
while (contador < 5) {
  Serial.println(contador);
  contador++;
}

// Aguardar dado serial
while (Serial.available() == 0) {
  delay(100);
}
char comando = Serial.read();
\end{lstlisting}
\end{frame}

\begin{frame}[fragile]{Inicializando array}
\begin{lstlisting}
// Forma 1: com valores explícitos
int valores[5] = {1, 2, 3, 4, 5};

// Forma 2: sem especificar tamanho
int valores[] = {10, 20, 30, 40};

// Forma 3: só declarar (valores são 0 ou lixo)
int sensores[8];

// Forma 4: iniciar todos com 0
int contador[10] = {0};
\end{lstlisting}
\end{frame}

\begin{frame}{Aplicação no seguidor}
  \begin{itemize}
    \item Use array para guardar leitura de múltiplos sensores.
    \item Processe cada sensor com laço.
    \item Calcule erro médio ou ponderado.
    \item Muito mais eficiente que variáveis individuais.
  \end{itemize}
\end{frame}

\begin{frame}[fragile]{Exemplo completo}
\begin{lstlisting}
const int NUM_SENSORES = 5;
int pinsSensores[] = {A0,A1,A2,A3,A4};
int leituras[NUM_SENSORES];

void setup() {
  for (int i = 0; i < NUM_SENSORES; i++) {
    pinMode(pinsSensores[i], INPUT);
  }
}

void loop() {
  for (int i = 0; i < NUM_SENSORES; i++) {
    leituras[i] = analogRead(pinsSensores[i]);
  }
  // Agora processe as leituras com PID...
}
\end{lstlisting}
\end{frame}

\end{document}
